\chapter{考试范围}


第1--6单元 + 名著《骆驼祥子》《钢铁是怎样炼成的》，覆盖基础、阅读、写作三大板块，按考查权重排序。附：必背要点与答题模板，适配所有部编版教材，七年级下册通用。

\section*{一、语言积累与运用（25--30分，基础必拿分）}

\subsection*{1. 字音字形（8--10分）}

\textbf{高频易错字：}

\begin{itemize}[leftmargin=*]
    \item \textbf{第一单元：} 殷（yān）红、卓（zhuó）越、鞠（jū）躬尽、锲（qiè）而不舍、锋芒毕露、至死不懈、家喻户晓
    \item \textbf{第二单元：} 哺（bǔ）育、澎湃（pài）、澜（lán）语、默契（qì）、浩浩荡荡
    \item \textbf{第三单元：} 憎（zēng）恶、区（ōu）楼、滞（zhì）留、颔（hàn）头、取缔（dì）、烦躁（zào）
    \item \textbf{第四单元：} 崎岖（qí qū）、噩耗（è hào）、毛坯（máo pī）、辗转难眠（zhǎn zhuǎn nán mián）
    \item \textbf{第五单元：} 深坳（shēn ào）、梭镖（suō biāo）、砍伐（kǎn fá）、摇撼（yáo hàn）
    \item \textbf{第六单元：} 合拢（hé lǒng）、吟唱（yín chàng）、孤零零（gū líng líng）、不期而至（bù qī ér zhì）、心有灵犀（xīn yǒu líng xī）
\end{itemize}

\textbf{多音字：} 调（diào/tiáo）、强（qiáng/qiǎng/jiàng）、塞（sāi/sè/sài）

\textbf{字形辨析：} 辨清“躁/燥/澡”“概/慨”“即/既”“辨/辩”，重点防范成语写错字。

\subsection*{2. 古诗文默写（6--8分）}

\textbf{必背篇目：} 《竹里馆》《春夜洛城闻笛》《晚春》《逢入京使》《孙权劝学》（重点句子）、《卖油翁》（重点句子）

\textbf{高频默写句（必背，杜绝错别字）：}

\begin{itemize}[leftmargin=*]
    \item \textbf{《竹里馆》：} 独坐幽篁里，弹琴复长啸。深林人不知，明月来相照。
    \item \textbf{《春夜洛城闻笛》：} 谁家玉笛暗飞声，散入春风满洛城。此夜曲中闻折柳，何人不起故园情。
    \item \textbf{《晚春》：} 草树知春不久归，百般红紫斗芳菲。杨花榆荚无才思，惟解漫天作雪飞。
    \item \textbf{《逢入京使》：} 故园东望路漫漫，双袖龙钟泪不干。马上相逢无纸笔，凭君传语报平安。
    \item \textbf{《孙权劝学》：} 卿今当涂掌事，不可不学！但当涉猎，见往事耳；士别三日，即更刮目相待。
\end{itemize}

\textbf{易错字：} 篁、折柳、涉猎、卿、涂、荚、惟、龙钟

\subsection*{3. 文学常识与名著（5--7分）}

\textbf{文学常识：}

\begin{itemize}[leftmargin=*]
    \item （1）\textbf{《资治通鉴》：} 北宋司马光主持编纂的编年体通史，记载了从战国到五代共1362年的历史。
    \item （2）\textbf{《木兰诗》：} 南北朝时期北方乐府民歌，与《孔雀东南飞》并称“乐府双璧”，收录于《乐府诗集》。
    \item （3）\textbf{重点作者及作品：} 杨振宁《邓稼先》、臧克家《说和做——记闻一多先生言行片段》、陆定一《老山界》、鲁迅《阿长与〈山海经〉》、欧阳修《卖油翁》。
\end{itemize}

\textbf{名著《骆驼祥子》（必考，重点掌握）：}

\begin{itemize}[leftmargin=*]
    \item （1）\textbf{作者：} 老舍，原名舒庆春，字舍予，现代著名作家，被授予“人民艺术家”称号。
    \item （2）\textbf{核心人物：} 祥子（前期：勤劳、朴实、坚韧、有梦想；后期：麻木、潦倒、自私、堕落）、虎妞（泼辣、霸道、贪婪）、刘四爷（残忍、自私、吝啬）。
    \item （3）\textbf{核心情节（三起三落，必考）：}
        \begin{itemize}
            \item 一起：祥子苦干三年，攒钱买了一辆新车；
            \item 一落：新车被大兵抢走，祥子沦为乞丐；
            \item 二起：祥子重新攒钱，准备买新车；
            \item 二落：攒下的钱被孙侦探敲诈一空；
            \item 三起：虎妞出钱，给祥子买了一辆二手车；
            \item 三落：虎妞难产而死，祥子卖掉二手车安葬虎妞，彻底沦为行尸走肉。
        \end{itemize}
    \item （4）\textbf{主题：} 揭露旧社会对底层劳动者的残酷压迫，批判黑暗、腐朽的社会制度，表达对底层人民的同情。
\end{itemize}

\textbf{名著《钢铁是怎样炼成的》（必考，重点掌握）：}

\begin{itemize}[leftmargin=*]
    \item （1）\textbf{作者：} 苏联奥斯特洛夫斯基，自传体长篇小说，全身瘫痪、双目失明后完成创作。
    \item （2）\textbf{核心人物：} 保尔（主角，意志坚定、顽强不屈、无私奉献，为人类解放事业奋斗终身）、朱赫来（保尔革命引路人，布尔什维克党员，教会保尔革命道理）。
    \item （3）\textbf{关键情节（填空/简答）：}
        \begin{itemize}
            \item 四次死里逃生：与波兰白军战斗，腿伤+伤寒；骑兵战头部重伤，右眼失明；肃反工作积劳成疾；筑路时患伤寒+大叶性肺炎。
        \end{itemize}
    \item （4）\textbf{保尔精神：} 为理想献身、坚韧不拔、顽强奋斗、无私奉献。
\end{itemize}

\subsection*{4. 词语运用（4--6分）}

\begin{itemize}[leftmargin=*]
    \item \textbf{成语辨析：} 重点掌握“锲而不舍、兀兀穷年、沥尽心血、家喻户晓、鲜为人知、可歌可泣、深恶痛绝、微不足道”，结合语境判断使用正误。
    \item \textbf{词语语境义：} 如“涉猎”（粗略阅读）、“刮目相待”（用新的眼光看待）、“微不足道”（非常渺小，不值得一提）。
\end{itemize}

\subsection*{5. 病句辨析与修改（3--4分）}

\textbf{高频病句类型：} 成分残缺、搭配不当、语序不当、前后矛盾、表意不明、重复累赘。

\textbf{典型示例及修改：}

\begin{enumerate}[leftmargin=*]
    \item 错误：通过开展“书香校园”活动，使同学们的阅读水平有了明显提高。\\
    修改：删去“通过”或“使”，补充主语。
    \item 错误：我们要养成认真完成作业的好习惯和态度。\\
    修改：删去“和态度”，搭配不当。
    \item 错误：他经常回忆过去的往事。\\
    修改：删去“过去的”或“往”，重复累赘。
\end{enumerate}

\section*{二、文言文阅读（12--15分，2篇：课内1篇+课外1篇）}

\subsection*{课内必考（考查概率90\%，重点突破）}

\subsubsection*{《短文两篇》——《陋室铭》}

\begin{table}[h]
\centering
\begin{tabular}{|m{4cm}|m{10cm}|}
\hline
\textbf{考查类型} & \textbf{具体内容} \\
\hline
词类活用 & 名、灵、馨、乱、劳、上 \\
\hline
古今异义 & 形、馨、鸿、丝竹、调（90\%学生混淆“鸿”的古今含义） \\
\hline
\multirow{5}{*}{重点句子翻译} & 1. 斯是陋室，惟吾德馨。 \\
 & 2. 苔痕上阶绿，草色入帘青。 \\
 & 3. 谈笑有鸿儒，往来无白丁。 \\
 & 4. 可以调素琴，阅金经。 \\
 & 5. 无丝竹之乱耳，无案牍之劳形。 \\
\hline
\end{tabular}
\end{table}